\section{研究方法论}

本部分描述所提出的研究设计，包括数据来源、样本构建、变量的操作化定义和分析策略。该方法论旨在实证检验第三节所规定的可检验假说，并在一个最大化内部有效性的框架内进行检验，同时充分承认来自无需许可金融协议观测数据的固有局限性。

\subsection{数据来源与协议选择}

本研究完全依赖来自基于以太坊的DeFi借贷协议的公开可用链上数据。数据来源的选择是有意为之，并且与研究设计的一个重要特征密切相关：由于以太坊区块链上的每一笔交易都是公开可见且可核实的，研究可由任何具备区块链基础设施访问权限的研究者独立复现，无需依赖专有数据集、平台合作协议或商业数据许可安排。

主要数据来源将是三大主流DeFi借贷协议的历史事件日志和状态数据：Aave（V2和V3）、Compound（V2和V3）以及MakerDAO。这三大协议共同占据了以太坊上DeFi借贷活动的大部分份额，并为跨协议验证提供了互补的数据。要提取的特定事件类型包括：（i）\texttt{Deposit}和\texttt{Withdraw}事件，代表抵押品的提供与提取；（ii）\texttt{Borrow}和\texttt{Repay}事件，代表借贷头寸的建立与偿还；（iii）\texttt{LiquidationCall}或等效事件，捕获第三方触发的清算执行；以及（iv）\texttt{FlashLoan}事件，用于识别可能涉及复杂多步策略而非简单借贷行为的交易记录。

数据将通过Dune Analytics提取，这是一种社区维护的区块链数据平台，提供对已解码事件日志和状态数据的SQL可查询访问。标的抵押资产价格数据将从链上价格预言机（Chainlink）获取，以确保分析中使用的价格信息与分析时刻协议参与者实际可获取的价格信息相匹配。研究期间将从各协议的创世开始跨越至分析时最近一个完整日历月的数据，确保涵盖多个市场周期和不同的DeFi生态活跃度水平。

\subsection{分析单元与样本构建}

主要的分析单元为\textbf{借款人-头寸-月份}，即由特定钱包地址在给定日历月内持有的特定借贷头寸。这一分析单元的选择基于以下考虑。第一，它使得能够追踪个体头寸随时间的演变，捕捉风险积累和贷方响应的动态过程——这正是本研究问题的核心所在。第二，它适应了单个钱包地址在协议内和跨协议可能同时持有多个独立借贷头寸的事实，并允许在钱包层面使用聚类标准误方法来控制误差的地址内相关性。第三，它与大多数DeFi风险指标常规报告和消费的时间频率保持一致。

样本将包含符合以下条件的所有借贷头寸：（i）借贷事件发生于协议在以太坊主网上部署期间或之后且早于数据提取日期；（ii）借贷头寸在观测窗口内的某一时间点具有非零健康因子；（iii）借贷头寸涉及协议支持的主要借贷资产（ETH、WBTC、USDC、USDT、DAI）之一。属于纯粹的杠杆化收益耕作策略，具有围绕同一资产的循环借贷模式（可识别）的头寸将被标记并与主样本分开分析，以评估结果的稳健性。

\subsection{主动与被动分类}

方法论的另一个关键要素是区分主动、借款人发起的操作与被动、外部驱动事件的可操作区分。此一区分至关重要，因为研究问题所涉及的是借款人的行为，而将主动借贷决策与第三方或协议机制触发的事件混为一谈将严重削弱分析的可证伪效力。

分类规则如下。当且仅当交易发送者（\texttt{msg.sender}）与借款人钱包地址一致时，一笔链上交易被分类为\textbf{主动借贷行为}。依据此一准则，以下事件分类为主动操作：（i）追加存入抵押品；（ii）偿还未偿还贷款；（iii）提取抵押品（在协议允许的情况下）；（iv）追加借款；以及（v）通过全额偿还未偿还债务并提取全部抵押品来平仓。当交易发送者非借款人地址（无论是已识别的清算机器人或是其他地址）时，该笔交易分类为\textbf{被动清算事件}。

该分类规则的重要优势在于其完全可决定性且可复现性：任何具备原始交易数据访问权限的研究者均可无歧义地应用该规则。其主要局限性在于无法区分借款人为交易的实际（真正主动行为）发起者，以及借款人将交易签名委托给第三方服务的情况，但分析可以通过限定样本至地址基于交易nonce模式判断大多数交易为自发起的地址来进行稳健性检验。

\subsection{头寸风险轨迹重构}

对于每一笔借贷头寸，将按每日频率重构头寸健康状况的变化轨迹。核心风险度量指标为\textbf{健康因子}（Health Factor, HF），其定义如下：

\begin{equation}
\text{HF}_t = \frac{\sum_{i} (C_{i,t} \cdot P_{i,t}) \cdot \text{LTV}_i}{D_t}
\end{equation}

其中\(C_{i,t}\)为时间\(t\)抵押资产\(i\)的数量，\(P_{i,t}\)为其价格，\(\text{LTV}_i\)为该资产的协议特定LTV参数，\(D_t\)为以通用货币单位计价的未偿还债务总额。HF在清算阈值处等于1.0，值越低表示风险越高。

轨迹重构还将计算一组刻画头寸时变演化过程的衍生变量：（i）过去7天、14天和30天内达到的最低健康因子；（ii）头寸在不同健康因子区间（如低于1.2、1.2--1.5、1.5--2.0、高于2.0）所经历的天数；（iii）过去30天内的最大健康因子回撤；以及（iv）头寸自我进入当前健康因子区间以来所经过的时间。

\subsection{行为过程变量}

主要关注的解释变量捕捉借款人对头寸风险变化的主动行为响应。对于每个借款人-头寸-月份，将构建以下变量：

\begin{description}[leftmargin=*,labelindent=\parindent, itemindent=0pt]
\item[主动抵押品调整率。] 借款人在观测月中主动追加或提取的抵押品净额，除以同期平均总抵押品价值。正值表示主动净追加抵押品（风险减轻），负值表示主动净提取（风险增加）。

\item[主动债务管理率。] 借款人在观测月中主动偿还或新增的债务净额，除以同期平均未偿还债务。正值表示新增净借款，负值表示净还款。

\item[响应延迟。] 从健康因子首次下穿给定阈值（如1.5、1.3或1.1）至借款人首次进行后续主动追加抵押品或偿还债务之间的天数。较长时延表明在风险上升面前的被动行为。

\item[调整强度。] 进入新风险区间后的首次调整中主动追加抵押品或偿还债务的绝对量级，以头寸总价值的比例表示。这衡量了借款人作出的是小幅尝试性调整还是大幅果断调整。

\item[行为一致性。] 度量借款人行为在时间维度相比于基于规则的调整策略的接近程度。健康因子下降时系统性追加抵押品的借款人将获得较高的一致性得分；行为看似随机或对价格波动而非对风险变化的响应者将获得较低得分。
\end{description}

\subsection{控制变量}

为分离出行为过程变量的增量贡献，模型将包含四个领域的综合控制变量集：

\begin{enumerate}[leftmargin=*,labelindent=\parindent, itemindent=0pt]
\item \textbf{当前风险状态：} 时点健康因子、抵押率以及波动性资产在总抵押品中的比例。
\item \textbf{账户特征：} 钱包年龄、跨所有DeFi协议的历史交易总次数、发生交互的不同协议数量以及历史借贷总价值。
\item \textbf{资产构成：} 具体用作抵押品和借入资产的哑变量，以及头寸内部资产集中度的赫芬达尔-赫希曼指数。
\item \textbf{市场条件：} 抵押资产过去30天的年化波动率、抵押资产与借入资产价格之间的相关性，以及可能抑制主动头寸管理的高gas费用期间哑变量。
\end{enumerate}

\subsection{分析框架}

分析将分三个阶段进行，与研究问题的层次结构相对应。

%
% Figure 2: Methodology Pipeline
%
\begin{figure}[t]
\centering
\resizebox{\textwidth}{!}{%
\begin{tikzpicture}[
  node distance=0.7cm and 0.4cm,
  stage/.style={rectangle, rounded corners=3pt, draw=#1!70, fill=#1!15, text=black!80, minimum width=2.8cm, minimum height=0.7cm, align=center, font=\scriptsize\sffamily, line width=0.4pt},
  substage/.style={rectangle, rounded corners=2pt, draw=#1!45, fill=#1!8, text=black!65, minimum width=2.6cm, minimum height=0.5cm, align=center, font=\tiny\sffamily, line width=0.3pt},
  arrow/.style={->, >=stealth, line width=0.6pt, draw=black!60},
]
  \node[stage=blue] (s1) {1. Data Extraction\\Ethereum RPC / Dune};
  \node[stage=cyan, right=0.08cm of s1] (s2) {2. Event Classification\\Active vs.\ Passive};
  \node[stage=teal, right=0.08cm of s2] (s3) {3. Trajectory\\Reconstruction\\Daily HF, bands};
  \node[stage=violet, right=0.08cm of s3] (s4) {4. Variable\\Construction\\BH process vars};
  \node[stage=red!70!black, right=0.08cm of s4] (s5) {5. Predictive\\Modeling\\(classification)};

  \draw[arrow, thick] (s1.east) -- (s2.west);
  \draw[arrow, thick] (s2.east) -- (s3.west);
  \draw[arrow, thick] (s3.east) -- (s4.west);
  \draw[arrow, thick] (s4.east) -- (s5.west);

  \node[substage=blue, below=0.1cm of s1] {\shortstack{Borrow, Repay,\\Deposit, Withdraw}};
  \node[substage=cyan, below=0.1cm of s2] {\shortstack{msg.sender\\classification}};
  \node[substage=teal, below=0.1cm of s3] {\shortstack{min HF, band\\days, drawdown}};
  \node[substage=violet, below=0.1cm of s4] {\shortstack{adjust rate,\\latency, intensity}};
  \node[substage=red!70!black, below=0.1cm of s5] {\shortstack{AUC-ROC, PR-AUC\\model comparison}};
\end{tikzpicture}%
}
\caption{Methodology pipeline from raw blockchain data to predictive model evaluation. Stage 1 extracts transaction events from Ethereum; Stage 2 classifies transactions as active (borrower-initiated) or passive (liquidation bot); Stage 3 reconstructs position risk trajectories; Stage 4 constructs behavioral process variables; and Stage 5 estimates predictive models using standard classification techniques, evaluating them $\rightarrow$ liquidation prediction at 7d/30d/90d horizons.}
\label{fig:methodology}
\end{figure}

\vspace{4pt}

\paragraph{第一阶段：行为描述（RQ1）。} 本阶段的目标是刻画头寸健康状况与主动借贷行为之间的经验关系。主要分析工具是估计如下面板回归模型：

\begin{equation}
\text{Action}_{i,t} = \alpha + \beta_1 \text{HF}_{i,t-1} + \beta_2 \text{HF}_{i,t-1}^2 + \gamma_1 \Delta P_{i,t} + \delta \mathbf{X}_{i,t-1} + \eta_i + \tau_t + \varepsilon_{i,t}
\end{equation}

其中\(\text{Action}_{i,t}\)为主动借贷行为（追加抵押品、偿还债务）的度量，\(\text{HF}_{i,t-1}\)为滞后健康因子，\(\Delta P_{i,t}\)为同一时期的价格变动，\(\mathbf{X}_{i,t-1}\)为时变控制变量向量，\(\eta_i\)通过固定效应捕获时变不可观测的异质性，\(\tau_t\)捕获共同时间效应。H1a通过\(\beta_1\)的符号和显著性进行检验，如果借款人通过采用风险减轻行动来应对健康因子的下降（即风险上升），则\(\beta_1\)应为负。H1b通过比较健康因子下降（损失域）和健康因子上升（收益域）的系数来检验。H1c通过增加头寸是否处于清算阈值附近窄带内的示性变量并将该示性变量与健康因子交互来扩展模型，从而实现检验。

\paragraph{第二阶段：预测框架（RQ2）。} 本阶段的目标是检验行为过程变量是否包含对清算事件的增量预测能力。分析框架为模型比较，在该逐步比较中将日益丰富的预测变量集与它们在保留测试样本中预测清算事件的能力进行比较。将估计以下模型：

\begin{itemize}[leftmargin=*,labelindent=\parindent, itemindent=0pt]
\item \textbf{M0（基准模型）：} 仅使用当前健康因子、抵押率及二者交互项的逻辑回归模型。
\item \textbf{M1（状态变量模型）：} M0在4.5节所描述的风险状态、账户特征和资产构成变量集的基础上增加变量的模型。
\item \textbf{M2（状态+市场变量模型）：} 在M1基础上增加市场条件变量的模型，包括抵押资产波动率和gas价格指标。
\item \textbf{M3（完全模型）：} 在M2基础上进一步增加第4.4节所描述行为过程变量的模型。
\end{itemize}

所有模型将主要使用梯度提升决策树（XGBoost）进行估计，并以逻辑回归结果作为稳健性报告。选择基于树的模型的原因在于其能够在无需研究者显式指定的情况下捕捉预测变量之间的非线性和复杂交互——在缺乏关于预测关系函数形式的强有力先验的条件下，这一点尤为重要。

模型性能将使用标准分类指标——ROC曲线下面积（AUC-ROC）、精确率-召回率AUC以及校正误差——在时间上分离的保留测试集上进行评估，采用模拟真实样本外条件的滚动窗口评估方案。模型间AUC差异的统计显著性将使用DeLong检验进行评定。若M3的AUC显著高于M2，则H2a获得支持。需要特别强调这一检验的严格性：由于M3包含M2中所有的变量再加上行为过程变量，任何观察到的预测性能改善均可直接归因于行为变量的信息含量，而不能归因于新增变量本身。

\paragraph{第三阶段：稳健性。} 一系列稳健性检验将评估结果对特定方法论选择的敏感性，包括：（i）使用逻辑回归而不是树模型重新估计所有模型，并比较树模型和线性规范的结果；（ii）改变预测期限（7天、30天或90天提前清算预测）；（iii）将样本限制在已在协议中至少交互最低天数的借款人，排除新用户；（iv）分别对每个协议重复分析以评估跨协议的可推广性；（v）对变量重要性进行分解，量化状态变量、市场变量和行为过程变量对总体预测能力的相对贡献。
